\chapter{Image Restoration} \label{ch:restoration}
\index{Image Restoration}

\section{Introduction to Image Denoising} \label{sec:intro_denoising}
\index{Image Denoising}

Image denoising is a fundamental inverse problem in computer vision and image processing. It plays a crucial role in multiple stages of imaging pipelines, serving as:
\begin{itemize}
    \item \textbf{Preprocessing}: to enhance the output quality and improve the efficacy of subsequent downstream algorithms.
    \item \textbf{Post-processing}: to remove compression artifacts, such as blocking and ringing effects.
    \item \textbf{Plug-and-play filters}: acting as an implicit regularization prior in various imaging applications.
\end{itemize}

The fundamental goal of an image denoising algorithm is to provide an estimate $\hat{y}$ of an unknown original image $y \colon X \to \mathbb{R}$ from an observed noisy image $z \colon X \to \mathbb{R}$. The observation model is typically formulated as:
\begin{equation}
    z(x) = y(x) + \eta(x), \quad x \in X
\end{equation}
where $\eta(x)$ represents independent and identically distributed (i.i.d.) white Gaussian noise, such that $\eta \sim \mathcal{N}(0, \sigma^2)$. While this formulation holds for grayscale images, it straightforwardly extends to RGB images.

\begin{insightbox}[Inverse Problems and Restoration Networks]
\index{Inverse Problems}
In the context of image restoration, the \textit{forward pass} is governed by well-defined physical laws and sensor characteristics that define how a clean signal degrades into a noisy observation. However, the \textit{inverse pass}---recovering the original signal from its degraded version---is inherently \textit{ill-posed}. Multiple clean images could plausibly map to the same noisy observation due to the loss of information caused by the perturbation. 

Deep restoration networks tackle this by learning a strong natural image prior from large datasets. Instead of analytically inverting the forward model, CNNs employ highly non-linear transformations to map the degraded space back to the manifold of natural images. By optimizing over vast amounts of data, the network essentially learns the statistical properties of clean images and the underlying noise distribution, effectively regularizing the ill-posed inverse problem.
\end{insightbox}

\section{Simple Deep Neural Networks for Image Denoising} \label{sec:dnn_denoising}

The advent of Convolutional Neural Networks (CNNs) has introduced a paradigm shift in image denoising, transitioning from traditional statistical modeling of input and noise to supervised learning approaches. 
Is denoising by deep neural networks a supervised or unsupervised learning problem? While the task resembles unsupervised learning since no explicit semantic labels are required, the training process is fundamentally supervised. The loss function is typically the Mean Squared Error (MSE), defined as $||\hat{y} - y||_2^2$, which assumes the availability of the noise-free ground truth image $y$. Alternatively, it can be viewed as self-supervised learning, as annotated pairs can be easily generated by adding synthetic noise to clean images, though the auxiliary task (denoising) is already the ultimate goal.

A straightforward solution involves designing a CNN with identical input and output spatial dimensions, akin to architectures used in semantic segmentation. The training data $TR = \{(y, z)\}$ consists of natural images $y$ and their noisy counterparts $z$. As long as the forward model is known and realistic, and the CNN can be sufficiently trained, this procedure can be generalized to any type of perturbation, including deblocking, deblurring, super-resolution, inpainting, enhancement, and de-raining.

\subsection{DnCNN: Residual Learning of Deep CNNs} \label{sec:dncnn}
\index{DnCNN} \index{Residual Learning}

A landmark architecture in this domain is the Denoising Convolutional Neural Network (DnCNN). Unlike traditional networks that aim to directly estimate the noise-free image $\hat{y}$, DnCNN is trained via \textit{residual learning} to predict the noise realization $\eta$. Given the relation $z = y + \eta$, the predicted noise-free image is recovered as:
\begin{equation}
    \hat{y} = z - \hat{\eta} = z - \text{CNN}(z)
\end{equation}
It is empirically easier to optimize the loss function over the noise residual $||\eta - \text{CNN}(z)||_2^2$ rather than the image directly. The network is fully convolutional (FCNN), allowing it to process images of arbitrary sizes.

The DnCNN architecture typically consists of:
\begin{itemize}
    \item \textbf{Layer 1}: Convolution + ReLU, featuring 64 filters of size $3 \times 3 \times c$.
    \item \textbf{Layers 2 to $d-1$}: Convolution + Batch Normalization (BN) + ReLU, with 64 filters of size $3 \times 3 \times 64$.
    \item \textbf{Layer $d$}: Convolution, using $c$ filters of size $3 \times 3 \times 64$ to reconstruct the input.
\end{itemize}
To mitigate border artifacts, zero padding is applied outside the image before processing. The loss function to minimize is formulated as:
\begin{equation}
    \ell(\Theta) = \frac{1}{2N} \sum_{i=1}^{N} ||\text{CNN}(z_i, \Theta) - (z_i - y_i)||_F^2
\end{equation}
where $||\cdot||_F$ denotes the Frobenius norm. Training is performed patch-wise (e.g., $40 \times 40$ patches for Gaussian denoising and $50 \times 50$ for others). Larger training sets beyond a certain threshold (e.g., 400 images of $180 \times 180$ pixels) do not substantially improve performance.

A significant limitation of DnCNN is that the network parameters $\Theta_\sigma$ are strongly dependent on the noise standard deviation $\sigma$. Consequently, a DnCNN trained on a specific noise level struggles to generalize to images corrupted by different noise levels, different noise models, or spatially variant (signal-dependent) noise. Utilizing a $\sigma$ higher than the true value results in over-smoothing (reducing PSNR), while a lower $\sigma$ yields noisy outputs.

\subsection{FFDNet: Fast and Flexible Denoising} \label{sec:ffdnet}
\index{FFDNet}

To overcome the rigidity of DnCNN, FFDNet introduces a flexible framework capable of handling spatially variant noise and multiple noise levels with a single network. FFDNet learns a mapping function $\mathcal{F}(\Theta, z, M)$, where the network parameters $\Theta$ are independent of the noise parameters. 

The inputs to FFDNet for an image of size $W \times H \times C$ include:
\begin{enumerate}
    \item A \textbf{noise level map} $M$ matching the spatial dimensions, indicating the noise standard deviation at each pixel.
    \item The input image $z$, which undergoes a \textbf{pixel shuffling} (downsampling) operation. The image is decomposed into 4 downsampled and cropped sub-images of size $\frac{W}{2} \times \frac{H}{2} \times C$, stacked across channels.
\end{enumerate}
The resulting input tensor has dimensions $\frac{W}{2} \times \frac{H}{2} \times (4C + 1)$. This downsampling reduces the spatial extent of the input, decreasing computational complexity in subsequent convolutional layers and effectively expanding the receptive field.

The noise level map $M$ dictates the noise standard deviation. For Additive White Gaussian Noise (AWGN), $M$ is a uniform image equal to $\sigma$. For spatially variant noise, $M$ is non-uniform, defined by the noisy image $z$ and sensor characteristics. The loss function minimized is:
\begin{equation}
    \ell(\Theta) = \frac{1}{2N} \sum_{i=1}^{N} ||\mathcal{F}(z_i, \Theta, M) - y_i||_F^2
\end{equation}

\section{Generalization and Bias-Free Networks} \label{sec:bias_free}
\index{Bias-Free Networks} \index{Generalization}

An important theoretical advantage of deep-learning techniques is the potential for a single neural network to denoise across a wide range of noise levels. However, empirical evidence shows that standard architectures systematically overfit to the specific noise levels present in their training set.

Generalization to unseen noise levels can be achieved through a simple, yet profound, architectural modification: the removal of all additive constants (biases) from the network. \textit{Bias-free} Convolutional Neural Networks can attain state-of-the-art performance across a broad spectrum of noise levels, even when trained exclusively on a narrow range of noise magnitudes. By enforcing a bias-free structure, the network relies strictly on linear combinations and multiplicative interactions, preventing the rigid memorization of specific intensity shifts associated with fixed noise levels.

\section{Class-Aware Denoising} \label{sec:class_aware}
\index{Class-Aware Denoising}

Standard denoising architectures process all images through a universal set of parameters. Class-aware denoising improves upon this by incorporating strong priors regarding the semantic content of the images. By narrowing down the manifold of valid images, a network trained on a specific class (e.g., faces, pets, flowers, living rooms, streets) can achieve superior restoration quality. 

During inference, a classification CNN first categorizes the input image, subsequently routing it to the appropriate class-specific fully convolutional denoiser. Key ingredients of class-aware networks include gradual denoising as the volume circulates through the layers, fast execution, and a relatively low parameter count. The architecture typically exploits residual connections where the noisy input is iteratively added to the gradual outputs, which act as noise estimates. Shallower layers handle fine details and high-frequency noise, while deeper layers, equipped with larger receptive fields, focus on recovering edges and structural textures distinguishing them from random noise.

\section{Self-Supervised Denoising} \label{sec:self_supervised}
\index{Self-Supervised Learning} \index{Noise2Noise} \index{Noise2Void} \index{Blind-Spot Networks}

Supervised denoising heavily relies on the availability of noisy-clean image pairs. However, acquiring clean ground truth images is often impractical in real-world scenarios, such as in medical imaging (CT scans, where higher doses of radiation are required for clean images) or outdoor photography in dynamic scenes. Self-supervised image denoising circumvents the need for clean images entirely.

\subsection{Noise2Noise (N2N)}
\index{Noise2Noise}
The Noise2Noise (N2N) framework demonstrates that it is possible to train a denoiser using only pairs of noisy images of the exact same underlying scene. Assuming the noise is zero-mean and independent between two noisy realizations $z = y + \eta$ and $z' = y + \eta'$, the optimization target can be shifted from the clean image $y$ to the alternate noisy realization $z'$:
\begin{equation}
    \Theta^* = \arg \min_{\Theta} \mathbb{E}_{z,z'} \left[ \mathcal{L}(\text{CNN}(\Theta; z), z') \right]
\end{equation}
On average, the optimal function learned by the network still predicts the clean signal $y$. N2N achieves performance comparable to supervised denoisers without explicit noise modeling or clean data. However, it still requires pairs of perfectly aligned noisy images, which can be challenging to obtain.

\subsection{Noise2Void (N2V) and Blind-Spot Networks}
\index{Noise2Void} \index{Blind-Spot Networks}
Noise2Void (N2V) further relaxes data requirements by learning from entirely \textit{unpaired} single noisy images. A naive approach of using the noisy image as both the input and the target would trivially lead the network to learn the identity function.

To prevent this, N2V introduces the \textbf{Blind-Spot Network (BSN)} training framework. The core idea is to hide specific pixels (blind spots) from the receptive field of the network. The network is then tasked to predict the value of the blind spot based solely on its neighborhood. This relies on the fundamental assumption that natural images possess strong spatial correlation, whereas the noise is pixel-wise independent and zero-mean.
The loss function focuses exclusively on the error at the blind spots:
\begin{equation}
    \Theta^* = \arg \min_{\Theta} \mathcal{L}(B .* \text{CNN}(\Theta; M(z)), B .* z)
\end{equation}
where $M$ is a masker function replacing the blind spot with a substitute value (e.g., a random neighbor, local mean/median, or uniform random value), and $B$ is a binary mask isolating the target pixels. Alternatively, blind spots can be architecturally enforced using specific structural constraints, such as half-plane convolutions or dilated convolutions with holes in the center of the filters.

\section{Blind Denoising of Real Photographs: CBDNet} \label{sec:cbdnet}
\index{CBDNet}

Standard CNN denoisers often overfit to synthetic Gaussian noise, generalizing poorly to real-world photographs that exhibit sophisticated, signal-dependent noise characteristics. Convolutional Blind Denoising (CBDNet) addresses this by incorporating a realistic noise model that accounts for the Poisson-Gaussian distribution inherent in photon counting, alongside the complete in-camera processing pipeline, including demosaicing, gamma correction, and JPEG compression.

CBDNet is structurally comprised of two subnetworks:
\begin{enumerate}
    \item A \textbf{noise estimation} subnetwork that predicts the noise level map.
    \item A \textbf{non-blind denoising} subnetwork that takes the noisy image and the estimated noise map to reconstruct the clean image.
\end{enumerate}
Crucially, CBDNet adopts an \textit{asymmetric loss function} that imposes a heavier penalty on the under-estimation of the noise level than on over-estimation, acknowledging that under-estimation leads to visually displeasing noise remnants, whereas slight over-estimation results in a more acceptable soft blurring effect. The overall architecture typically includes a bottleneck to consolidate representations before final reconstruction.

\section{Concluding Remarks} \label{sec:conclusion_restoration}
\index{Image Restoration!Conclusion}

The domain of image restoration has witnessed a vivid resurgence driven by deep learning. While traditional algorithms initially held efficiency advantages, deep CNNs now firmly dictate the state-of-the-art in performance. Contemporary research is heavily influenced by classical priors---such as self-similarity, sparsity, and multiscale processing---which are now embedded into complex neural architectures and iterative unfolding schemes.

Current research frontiers focus on developing highly efficient architectures for large images, increasing robustness against complex real-world noise models, and pushing the boundaries of practical training paradigms via self-supervised and blind-spot methodologies that eliminate the reliance on idealized ground truth data.
